\Needspace{12\baselineskip}
\originalchapter{热输运、浮力与色散}
\originalsection{热量守恒和Boussinesq模型}
常热容、导热系数流体满足 $\rho c_pDT/Dt=k_T\nabla^2T+\Phi+Q$。低速、温差较小时可以忽略黏性热和某些压缩效应，但须先比较尺度。Boussinesq近似把密度变化只保留在重力项中：$\rho=\rho_0[1-\alpha(T-T_0)]$，速度仍满足不可压缩约束，浮力加速度 $g\alpha(T-T_0)\hat z$。

下热上冷的平层静止传导态温度线性随高度变化。流体若略上移，较暖扰动产生浮力继续推动上升，同时黏性和热扩散会消除扰动。相应竞争量为 $Ra=g\alpha\Delta T H^3/(\nu\kappa)$，$Pr=\nu/\kappa$。临界值依赖热边界和机械边界；自由表面、刚性面和多孔层不是同一个数字。

\originalsection{羽流与分层}
局部热源驱动向上的羽流，卷吸周围流体使体积通量增长，浮力通量在无热损失时可以保持。层宽、中心速度和温差由动量、热量和卷吸或黏性假设共同决定。一个温度差不是独立的驱动力常数，它沿上升过程可能被稀释。

稳定分层中 $N_b^2=-(g/\rho_0)\dd\bar\rho/\dd z>0$，被抬高的流体较周围更重，产生恢复力。忽略阻尼的局部竖向扰动满足 $\ddot\zeta+N_b^2\zeta=0$，这联系了山后波和内部重力波。选择性抽取会涉及流线变形、扩散和稳定层的进入厚度，不能仅靠这个局部振子给出完整抽水效率。

\originalsection{剪切增强色散}
示踪浓度 $c$ 受对流扩散 $c_t+u(y)c_x=D(c_{xx}+c_{yy})$。横向扩散使粒子在快慢流层间交换，长期轴向扩散会比分子值大。取截面平均 $C(x,t)$，令 $c=C+\chi(y)C_x+\cdots$，在随均速 $\bar U$ 的长尺度中，横向胞问题为 $D\chi''=u-\bar U$，边界无通量、$\avg\chi=0$。

有效扩散 $D_{\rm eff}=D-\avg{(u-\bar U)\chi}=D+D\avg{(\chi')^2}$。第二等式来自分部积分及边界条件，明确给出增强量非负。它适用于时间远长于横向混合时间 $h^2/D$、轴向变化足够缓慢的条件；刚注入时的非高斯展宽不必服从同一长期扩散方程。

圆管Poiseuille流给经典增强 $\bar U^2a^2/(48D)$，开放明渠的边界和剖面不同，系数不能照搬。所选原题中实际解胞问题得到 $8U_0^2h^2/(945D)$，并明确 $U_0$ 的表面速度约定。

\originalsection{内能与焓：Boussinesq近似删去什么}
一般内能守恒为 $\rho De/Dt=-p\nabla\cdot\mathbf u+\tau_{ij}\partial_j u_i-\nabla\cdot\mathbf q+\rho r$，热流 $\mathbf q=-k_T\nabla T$。定义焓 $h=e+p/\rho$，结合质量方程有 $\rho Dh/Dt=Dp/Dt+\Phi-\nabla\cdot\mathbf q+\rho r$。因此压缩功并非在所有热问题中消失；常热容近不可压模型还需证明压力热效应相对导热/对流小。

若 $\alpha\Delta T\ll1$、流速远低于声速，声波快于关注的热流时间，可令密度变化仅进入浮力而用 $\nabla\cdot\mathbf u=0$。将静水压力从总压力中分离，扰动压力负责保证不可压，浮力 $b=-g\rho'/\rho_0$。浮力项可能虽相对总重力小，却与其他驱动项同阶，所以不能因为密度变化小就全部删去。

\originalsection{分层内波的色散、能量与传播}
线性稳定分层采用 $u_t=-p_x/\rho_0,w_t=-p_z/\rho_0+b,b_t+N_b^2w=0,u_x+w_z=0$。平面波 $\e^{i(kx+mz-\omega t)}$ 消去 $u,p,b$ 得
\[
\omega^2=N_b^2\frac{k^2}{k^2+m^2}.
\]
频率只取决于波矢方向，$0<|\omega|\le N_b$。取正频分支，群速 $\partial\omega/\partial(k,m)$ 与波矢垂直，所以相位与能量传播方向不一致。将山后波的条纹方向直接当成能量通量方向会错。

扰动能密度为 $\rho_0[u^2+w^2+b^2/N_b^2]/2$。前两项为动能，最后一项来自浮力势能；用线性方程相乘相加得 $\partial_t E+\nabla\cdot(p'\mathbf u)=0$。黏性/热扩散引入负耗散，稳定层仍可传送波能，不能因为流体“回到原位”就说没有输运。

定常地形、背景风 $U$ 中本征频率为 $\omega-kU=-kU$。内波关系给 $m^2=N_b^2/U^2-k^2$，只有 $|kU|<N_b$ 才有传播垂向波数；长波近静水时 $|m|\simeq N_b/U$。选择向上传能的分支需检查群速度，而非任意取 $m>0$。

\originalsection{坡面热风的耦合解}
考虑坡面与水平夹角 $\alpha_s$，沿坡速度 $u(n)$、法向距 $n$，背景垂向稳定温度梯度 $G>0$。小扰动温度 $\vartheta$ 的稳态近似给
\[
\nu u''+g\alpha_T\vartheta\sin\alpha_s=0,\qquad
\kappa\vartheta''=G u\sin\alpha_s.
\]
消去温度有 $u''''+[N_b^2\sin^2\alpha_s/(\nu\kappa)]u=0$，$N_b^2=g\alpha_TG$。无滑移、壁温扰动 $\vartheta(0)=\vartheta_0$、远场衰减，解为
\[
u=A\e^{-n/\delta}\sin(n/\delta),\qquad
\vartheta=\vartheta_0\e^{-n/\delta}\cos(n/\delta),
\]
\[
\delta=\left(\frac{4\nu\kappa}{N_b^2\sin^2\alpha_s}\right)^{1/4},\qquad
A=\frac{g\alpha_T\vartheta_0\sin\alpha_s\,\delta^2}{2\nu}.
\]
墙面加热使沿坡分量上升，冷却改变符号。相位错位与法向反向小流来自动量/热量两个扩散过程耦合，不能仅用“热空气上升”解释剖面。该解假设平面坡、线性背景和常系数，实际日夜变化需非稳态边界。

\originalsection{层流热羽流的相似尺度}
二维上升热羽流，固定线热源，远处沿程热/浮力通量 $B=\int wb\dd x$ 守恒。边界层动量尺度 $w^2/z\sim b$ 与横向黏性 $\nu w/\delta^2$ 同阶；热扩散给 $w/z\sim\kappa/\delta^2$，固定 $Pr$ 时有相同尺度结构。于是 $\delta^2\sim\nu z/w$，再以 $w b\delta\sim B$ 得
\[
w_c\propto z^{1/5},\quad\delta\propto z^{2/5},\quad b_c\propto z^{-3/5}.
\]
动量会因持续浮力输入增长，所以不能像无浮力射流那样设动量通量常数。相似系数和剖面由耦合ODE与源通量归一化给出，依赖Pr数。

作为对比，湍流轴对称羽流若用固定卷吸比例、守恒浮力通量，可得到 $\delta\propto z,w_c\propto B^{1/3}z^{-1/3},Q\propto B^{1/3}z^{5/3}$。这是不同闭合假设的尺度说明，不把它冒作前述层流解或源课程的实测曲线。

\originalsection{分层水库的选择性抽取}
抽水希望取某一温度/密度层，但速度使邻层界面变形。无扩散线性稳定分层下，抽取区厚度由惯性时间与浮力恢复时间竞争。点汇流量 $Q$ 的三维速度尺度 $U\sim Q/\delta^2$，$U/\delta\sim N_b$ 给 $\delta\sim(Q/N_b)^{1/3}$；二维线汇的单位宽度流量 $q$ 给 $U\sim q/\delta$，因此 $\delta\sim(q/N_b)^{1/2}$。维数和流量单位不同，两个指数不能混用。

若抽取主要由黏性/热扩散控制，必须求温度扰动与速度的耦合边值问题，参数还含 $\nu,\kappa$。源讲义展示该类分层抽取与边界条件；以上惯性估计只在相应尺度区间有效，不代替所有热边界的精确效率。选择性不能只按出口管径判断，还需知道流量、背景梯度、混合与抽取持续时间。

\originalsection{Taylor色散的尺度消去与振荡推广}
令 $t\gg h^2/D$，截面浓度近均匀，缓慢轴向梯度激发小修正 $\chi(y)C_x$。横向胞方程 $D\chi''=u-\bar U$ 由快方向扩散平衡产生；条件 $\avg\chi=0$ 保证 $C$ 仍是平均浓度。修正带来平均轴向通量 $\avg{uc}=\bar UC+\avg{(u-\bar U)\chi}C_x$，于是有效扩散与正文式一致。

振荡流 $u(y,t)=\Re[\hat u(y)\e^{i\omega t}]$ 需时间周期胞解，$(D\partial_y^2-i\omega)\hat\chi=\hat u-\avg{\hat u}$，无通量边界。周期平均增强
\[
D_{\rm eff}-D=-\tfrac12\Re\avg{(\hat u-\avg{\hat u})\hat\chi^*}
=\frac D2\avg{|\hat\chi'|^2}\ge0.
\]
高频下粒子来不及充分跨快慢层，增强受到抑制；不能将瞬时速度平方平均后总套稳态系数。不同剖面和频率应实际解胞问题，源课的定常/周期色散因此共用一条推导主线。
